% Section 1.1: Vector spaces, bases and change of basis | Part: 1 | Depends on: none | Status: done | Original lines: 353-397
%!TEX root = ../../main.tex
\subsection{Vector spaces, bases and change of basis}\label{sec:vs}


\begin{formal}[{Vector space, basis, components}]
% REWRITTEN: notation E -> V, m -> n (docs/NOTATION.md)
A \emph{real vector space} $V$ is a set with addition and multiplication by real numbers satisfying the usual axioms. A \emph{basis} $\beta=\{e_1,\dots,e_n\}$ is a set of vectors such that every $v\in V$ can be written \emph{uniquely} as
\[
v=v^1e_1+\dots+v^ne_n = v^ie_i .
\]
The numbers $(v^1,\dots,v^n)\in\R^n$ are the \emph{components} (or coordinates) of $v$ with respect to $\beta$. Choosing a basis gives an isomorphism $V\to\R^n$, $v\mapsto (v^i)$.
\end{formal}

\begin{intuition}[{a vector is an arrow, components are a description of it}]
A vector is a geometric object (``an arrow'', a velocity, a displacement). It exists \emph{before} we choose a basis. The list of numbers $(v^i)$ is only a \emph{description} of the arrow relative to the rulers $e_i$ we decided to use. Change the rulers, and the numbers change, but the arrow does not. The whole of tensor calculus is a machine to keep track of which statements depend on the description and which do not.
\end{intuition}

\subsection*{Change of basis}
Let $\beta=\{e_i\}$ be the ``old'' basis and $\beta'=\{e'_i\}$ the ``new'' one. Write the new basis vectors in terms of the old ones:
\begin{equation}\label{eq:basischange}
e'_i=\sum_j A^j{}_i\, e_j ,\qquad A=(A^j{}_i)\in GL_n(\R)
\end{equation}
(the coefficients of $e'_i$ form the $i$-th \emph{column} of $A$). If $v=v^je_j=v'^ie'_i$, substituting \eqref{eq:basischange} gives
\begin{equation}\label{eq:compchange}
v^j=A^j{}_i\,v'^i\qquad\Longleftrightarrow\qquad v'^i=(A^{-1})^i{}_j\,v^j .
\end{equation}

\begin{intuition}[{basis vectors and components transform in opposite ways}]
Equation \eqref{eq:basischange} uses $A$ (``forward''), equation \eqref{eq:compchange} uses $A^{-1}$ (``backward''). This is natural: if you double the length of your ruler (basis vector), you need \emph{half} as many rulers to describe the same arrow. Objects whose components transform like $A^{-1}$ (opposite to the basis) are called \textbf{contravariant}; vectors are the prime example.
\end{intuition}

\begin{example}[{a change of basis in $\R^2$}]
Let $e_1=(1,0)$, $e_2=(0,1)$ and $e'_1=(1,1)$, $e'_2=(-1,1)$. Then
\[
A=\begin{pmatrix}1&-1\\1&1\end{pmatrix},\qquad A^{-1}=\tfrac12\begin{pmatrix}1&1\\-1&1\end{pmatrix}.
\]
The vector $v=(3,1)$ has old components $v=(3,1)$ and new components
\[
v'=A^{-1}v=\tfrac12\begin{pmatrix}4\\-2\end{pmatrix}=\begin{pmatrix}2\\-1\end{pmatrix}.
\]
Check: $2e'_1-e'_2=(2,2)-(-1,1)=(3,1)=v$. \checkmark\ The \emph{same arrow} has two different descriptions.
\end{example}
